\documentclass[12pt]{article}
\usepackage[utf8]{inputenc}
\usepackage[T1]{fontenc}
\usepackage{hyperref}
\title{Rutul Language Resources for K-12 Curriculum Development}
\author{Compiled by Claude Sonnet 5.5 (claude-sonnet-5-5) for Andr\'as Kornai}
\date{}
\begin{document}
\maketitle
\section{Introduction}
Rutul (ISO 639-3 rut) is a Nakh-Daghestanian language of Dagestan (Russia) and Azerbaijan. This dossier comes from one search. Pages were not opened, and publication details are unverified.

\section{Primary Educational Resources}
A Wikipedia mirror says that a first-grade ABC book appeared in 1992, followed by further school textbooks. No textbook URL was found. \url{https://en.wikipedia.org/wiki/Rutul_language}

\section{Linguistic and Reference Materials}
\begin{itemize}
\item Ismailova, article on compiling a Russian-Rutul terminological dictionary. \url{https://gramota.net/eng/materials/2/2016/8-1/34.html}
\item Ismailova Eminat Ibragimovna, a second Gramota.net item. \url{https://gramota.net/eng/materials/2/2017/7-1/25.html}
\item Pedagogy-related article on Gramota.net. \url{https://www.gramota.net/en/article/phil20161834/fulltext}
\item Online Kina Rutul dictionary on linghub.ru, built from field recordings made in 2016 to 2018. It is incomplete. \url{https://linghub.ru/rutuldict/}
\item Freelang Rutul dictionary. \url{https://www.freelang.net/dictionary/rutul.php}
\end{itemize}
The Wikipedia mirror also cites Rutul-Russian dictionaries from 2006 (Dzhamalov and Semedov) and 2019 (Alisultanov and Suleimanova). No URLs for them were found.

\section{Religious and Community Resources}
None found.

\section{Digital and Audio Resources}
The linghub.ru dictionary is the only digital resource found.

\section{Curriculum Development Context}
Rutul has printed school textbooks according to secondary sources. They need to be located and checked through library catalogues.
\end{document}